\textbf{Phần 1. Trắc nghiệm lựa chọn}
\setcounterpageref{ex}{0}

\begin{ex} %[3H3V79]
Một hình chữ nhật có chiều dài $24\text{ cm}$, chiều rộng $8\text{ cm}$. Chu vi hình chữ nhật đó là
\choice
{\True $64\text{ cm}$}
{$32\text{ cm}$}
{$192\text{ cm}$}
{$48\text{ cm}$}
\loigiai{
Chu vi: $(24 + 8) \times 2 = 64\text{ (cm)}$.
\textbf{Đáp án của câu hỏi là A}
}
\end{ex}

\begin{ex} %[3H3V79]
Một hình vuông có chu vi $48\text{ cm}$. Diện tích của hình vuông đó là
\choice
{$96\text{ cm}^2$}
{$128\text{ cm}^2$}
{\True $144\text{ cm}^2$}
{$120\text{ cm}^2$}
\loigiai{
Cạnh: $48 : 4 = 12\text{ cm}$. Diện tích: $12 \times 12 = 144\text{ (cm}^2)$.
\textbf{Đáp án của câu hỏi là C}
}
\end{ex}

\begin{ex} %[3H3V79]
Một hình tròn có đường kính $26\text{ cm}$. Bán kính của hình tròn đó là
\choice
{$52\text{ cm}$}
{\True $13\text{ cm}$}
{$26\text{ cm}$}
{$18\text{ cm}$}
\loigiai{
Bán kính: $26 : 2 = 13\text{ (cm)}$.
\textbf{Đáp án của câu hỏi là B}
}
\end{ex}

\begin{ex} %[3D3V79]
Đổi $3\text{ m } 5\text{ mm}$ ra đơn vị mi-li-mét:
\choice
{$350\text{ mm}$}
{$3050\text{ mm}$}
{$3500\text{ mm}$}
{\True $3005\text{ mm}$}
\loigiai{
$3\text{ m} = 3000\text{ mm}$. Tổng: $3000 + 5 = 3005\text{ mm}$.
\textbf{Đáp án của câu hỏi là D}
}
\end{ex}

\begin{ex} %[3D3V79]
Đổi $2\text{ kg } 80\text{ g}$ ra đơn vị gam:
\choice
{\True $2080\text{ g}$}
{$2800\text{ g}$}
{$280\text{ g}$}
{$2008\text{ g}$}
\loigiai{
$2\text{ kg} = 2000\text{ g}$. Tổng: $2000 + 80 = 2080\text{ g}$.
\textbf{Đáp án của câu hỏi là A}
}
\end{ex}

\begin{ex} %[3D3V79]
Một bình chứa $2\ell$ nước. Rót ra 5 cốc nước, mỗi cốc $250\text{ ml}$. Trong bình còn lại bao nhiêu mi-li-lít nước?
\choice
{$1250\text{ ml}$}
{$500\text{ ml}$}
{\True $750\text{ ml}$}
{$1000\text{ ml}$}
\loigiai{
Đổi $2\ell = 2000\text{ ml}$. Đã rót: $250 \times 5 = 1250\text{ ml}$. Còn lại: $2000 - 1250 = 750\text{ (ml)}$.
\textbf{Đáp án của câu hỏi là C}
}
\end{ex}

\begin{ex} %[3H3V79]
Một khối lập phương có bao nhiêu đỉnh, bao nhiêu cạnh và bao nhiêu mặt?
\choice
{$6$ đỉnh, $12$ cạnh, $8$ mặt}
{\True $8$ đỉnh, $12$ cạnh, $6$ mặt}
{$8$ đỉnh, $6$ cạnh, $12$ mặt}
{$12$ đỉnh, $8$ cạnh, $6$ mặt}
\loigiai{
Khối lập phương có 8 đỉnh, 12 cạnh, 6 mặt.
\textbf{Đáp án của câu hỏi là B}
}
\end{ex}

\begin{ex} %[3D3V79]
Nhiệt độ nước đá đang tan là
\choice
{$100^\circ\text{C}$}
{$37^\circ\text{C}$}
{$10^\circ\text{C}$}
{\True $0^\circ\text{C}$}
\loigiai{
Nước đá tan ở $0^\circ\text{C}$.
\textbf{Đáp án của câu hỏi là D}
}
\end{ex}

\begin{ex} %[3H3V79]
Cho đoạn thẳng $AB = 84\text{ cm}$. Gọi $M$ là trung điểm của $AB$, $N$ là trung điểm của $AM$. Độ dài đoạn thẳng $NB$ là
\choice
{\True $63\text{ cm}$}
{$42\text{ cm}$}
{$21\text{ cm}$}
{$54\text{ cm}$}
\loigiai{
$AM = 42\text{ cm}$; $AN = 21\text{ cm}$; $NB = 84 - 21 = 63\text{ (cm)}$.
\textbf{Đáp án của câu hỏi là A}
}
\end{ex}

\begin{ex} %[3H3V79]
Góc tạo bởi kim giờ và kim phút của đồng hồ lúc đúng 9 giờ là
\choice
{Góc nhọn}
{Góc tù}
{\True Góc vuông}
{Góc bẹt}
\loigiai{
Lúc 9 giờ đúng hai kim vuông góc với nhau.
\textbf{Đáp án của câu hỏi là C}
}
\end{ex}

\begin{ex} %[3H3V79]
Một hình chữ nhật có diện tích $80\text{ cm}^2$, chiều rộng $8\text{ cm}$. Chu vi hình chữ nhật đó là
\choice
{$40\text{ cm}$}
{\True $36\text{ cm}$}
{$32\text{ cm}$}
{$28\text{ cm}$}
\loigiai{
Chiều dài: $80 : 8 = 10\text{ cm}$. Chu vi: $(10 + 8) \times 2 = 36\text{ (cm)}$.
\textbf{Đáp án của câu hỏi là B}
}
\end{ex}

\begin{ex} %[3H3V79]
Cần bao nhiêu khối lập phương nhỏ cạnh $1\text{ cm}$ để xếp thành khối lập phương lớn cạnh $4\text{ cm}$?
\choice
{$16$ khối}
{$32$ khối}
{$48$ khối}
{\True $64$ khối}
\loigiai{
$4 \times 4 \times 4 = 64$ khối.
\textbf{Đáp án của câu hỏi là D}
}
\end{ex}

\textbf{Phần 2. Trắc nghiệm đúng sai}
\setcounterpageref{ex}{0}

\begin{ex} %[3H3V79]
Xét tính đúng sai của các kiến thức hình học và đo lường:
\choiceTF
{\True Hình vuông có 4 cạnh bằng nhau và 4 góc vuông.}
{$1\text{ kg} = 100\text{ g}$.}
{\True Đường kính của hình tròn gấp đôi bán kính.}
{Diện tích hình chữ nhật bằng (dài + rộng) nhân 2.}
\loigiai{
\begin{itemize}
    \item Ý a) Đúng: Định nghĩa chuẩn.
    \item Ý b) Sai: Sai vì $1\text{ kg} = 1000\text{ g}$.
    \item Ý c) Đúng: Mối quan hệ chuẩn.
    \item Ý d) Sai: Sai vì đó là chu vi.
\end{itemize}
\begin{itemize}
    \item a) \textbf{Đúng}.
    \item b) \textbf{Sai}.
    \item c) \textbf{Đúng}.
    \item d) \textbf{Sai}.
\end{itemize}
}
\end{ex}

\begin{ex} %[3H3V79]
Một khu đất hình chữ nhật có chiều dài $50\text{ m}$, chiều rộng $30\text{ m}$:
\choiceTF
{\True Chu vi khu đất là $160\text{ m}$.}
{\True Diện tích khu đất là $1500\text{ m}^2$.}
{Nếu bớt chiều dài đi $20\text{ m}$ thì khu đất có diện tích không đổi.}
{\True Xây hàng rào xung quanh khu đất chừa lối đi $4\text{ m}$ cần $156\text{ m}$ rào.}
\loigiai{
\begin{itemize}
    \item Ý a) Đúng: $(50 + 30) \times 2 = 160\text{ m}$.
    \item Ý b) Đúng: $50 \times 30 = 1500\text{ m}^2$.
    \item Ý c) Sai: Sai vì diện tích giảm xuống còn $30 \times 30 = 900\text{ m}^2$.
    \item Ý d) Đúng: $160 - 4 = 156\text{ m}$.
\end{itemize}
\begin{itemize}
    \item a) \textbf{Đúng}.
    \item b) \textbf{Đúng}.
    \item c) \textbf{Sai}.
    \item d) \textbf{Đúng}.
\end{itemize}
}
\end{ex}

\textbf{Phần 3. Trả lời ngắn}
\setcounterpageref{ex}{0}

\begin{ex} %[3H3V79]
Một hình vuông có chu vi $60\text{ cm}$. Tính diện tích hình vuông theo xăng-ti-mét vuông.
\shortans{225}
\loigiai{
\textbf{Đáp số:} $225$\\
Cạnh: $60 : 4 = 15\text{ cm}$. Diện tích: $15 \times 15 = 225\text{ (cm}^2)$.
}
\end{ex}

\begin{ex} %[3H3V79]
Hình tròn có bán kính $16\text{ cm}$. Tính đường kính hình tròn theo xăng-ti-mét.
\shortans{32}
\loigiai{
\textbf{Đáp số:} $32$\\
$16 \times 2 = 32\text{ (cm)}$.
}
\end{ex}

\begin{ex} %[3D3V79]
Đổi $5\text{ m } 8\text{ cm}$ ra mi-li-mét.
\shortans{5080}
\loigiai{
\textbf{Đáp số:} $5080$\\
$5000 + 80 = 5080\text{ mm}$.
}
\end{ex}

\begin{ex} %[3D3V79]
Một túi muối nặng $1\text{ kg}$. Dùng hết $350\text{ g}$. Túi muối còn lại bao nhiêu gam?
\shortans{650}
\loigiai{
\textbf{Đáp số:} $650$\\
$1000 - 350 = 650\text{ (g)}$.
}
\end{ex}

\begin{ex} %[3D3V79]
Đổi $3\ell - 850\text{ ml}$ ra mi-li-lít.
\shortans{2150}
\loigiai{
\textbf{Đáp số:} $2150$\\
$3000 - 850 = 2150\text{ (ml)}$.
}
\end{ex}

\begin{ex} %[3H3V79]
Một khối lập phương có chu vi 1 mặt là $36\text{ cm}$. Tính tổng độ dài 12 cạnh của khối lập phương theo xăng-ti-mét.
\shortans{108}
\loigiai{
\textbf{Đáp số:} $108$\\
Cạnh: $36 : 4 = 9\text{ cm}$. Tổng 12 cạnh: $12 \times 9 = 108\text{ (cm)}$.
}
\end{ex}
